% =============================================================================
% UFFIZI RL: THE 15-MINUTE TALK  ("The Art of Queueing")
% Fuller slides: \small body, larger figures, more content, distributed so each
% slide fills the vertical space.
% Build: pdflatex uffizi_talk.tex (twice for the frame count).
% =============================================================================
\documentclass[aspectratio=169, 11pt, t]{beamer}

\usetheme{Madrid}
\usecolortheme{whale}
\setbeamertemplate{navigation symbols}{}
\setbeamertemplate{footline}{%
  \leavevmode\hbox{%
    \begin{beamercolorbox}[wd=.5\paperwidth,ht=2.5ex,dp=1ex,center]{author in head/foot}%
      \usebeamerfont{author in head/foot}\insertshortauthor
    \end{beamercolorbox}%
    \begin{beamercolorbox}[wd=.5\paperwidth,ht=2.5ex,dp=1ex,right]{date in head/foot}%
      \usebeamerfont{date in head/foot}\insertframenumber{} / \inserttotalframenumber\hspace*{2ex}
    \end{beamercolorbox}}%
  \vskip0pt%
}
\setbeamerfont{frametitle}{size=\large,series=\bfseries}

\usepackage[T1]{fontenc}
\usepackage{lmodern}
\usepackage{amsmath,amssymb}
\usepackage{booktabs}
\usepackage{graphicx}
\usepackage{xcolor}

\definecolor{safegreen}{RGB}{60,160,60}
\definecolor{warnorange}{RGB}{200,120,20}
\definecolor{crowdred}{RGB}{200,50,50}
\newcommand{\good}[1]{\textcolor{safegreen}{#1}}
\newcommand{\warn}[1]{\textcolor{warnorange}{#1}}
\newcommand{\bad}[1]{\textcolor{crowdred}{#1}}
\newcommand{\emp}[1]{\textit{#1}}
\newcommand{\code}[1]{\texttt{\small #1}}
\newcommand{\fig}[1]{../outputs/figures/#1}
\newcommand{\asset}[1]{../outputs/assets/#1}
\newcommand{\pic}[1]{../outputs/assets/masterpieces/#1}

% Body font for content frames and roomy paragraph spacing, so slides fill.
\newcommand{\body}{\small}
\setlength{\parskip}{0.45em}

% Symptom / Diagnosis / Fix / Lesson, spread by \vfill so it fills top-to-bottom.
\newcommand{\sdfl}[4]{%
{\bad{\textbf{S}}\ \emp{Symptom.}\ #1}\par\vfill
{\warn{\textbf{D}}\ \emp{Diagnosis.}\ #2}\par\vfill
{\good{\textbf{F}}\ \emp{Fix.}\ #3}\par\vfill
{\textbf{L}\ \emp{Lesson.}\ #4}\par}

\title[The Art of Queueing]{The Art of Queueing}
\subtitle{Reinforcement Learning for visitor flow at the Gallerie degli Uffizi in Florence}
\author[Canal / Duarte Gon\c{c}alves / Guinea]{Marco Canal \and Michael Duarte Gon\c{c}alves \and Mircea-Adrian Guinea}
\institute[BSE]{Reinforcement Learning, Barcelona School of Economics}
\date{June 2026}

\begin{document}

% --- 1. Title ---------------------------------------------------------------
{\setbeamertemplate{footline}{}\begin{frame}[c]\titlepage\end{frame}}

% --- 2. The masterpieces ----------------------------------------------------
\begin{frame}{The rooms everyone queues for}
\body
Almost every one of the $\sim$5{,}000 daily visitors comes for the same handful of paintings,
crammed into just three rooms of the ninety-eight.
\vfill
\centering
\includegraphics[height=2.3cm]{\pic{botticelli_primavera.jpg}}\hfill
\includegraphics[height=2.3cm]{\pic{botticelli_venus.jpg}}\hfill
\includegraphics[height=2.3cm]{\pic{raphael_cardellino.jpg}}\\[10pt]
\includegraphics[height=2.3cm]{\pic{leonardo_annunciation.jpg}}\hfill
\includegraphics[height=2.3cm]{\pic{leonardo_magi.jpg}}\hfill
\includegraphics[height=2.3cm]{\pic{michelangelo_tondo.png}}
\vfill
\raggedright
Botticelli's \emp{Primavera} and \emp{Birth of Venus} (room A11/A12); Leonardo's
\emp{Annunciation} and \emp{Adoration of the Magi} (A35); Raphael's \emp{Madonna del Cardellino}
with Michelangelo's \emp{Doni Tondo} (A38).
\end{frame}

% --- 3. The two floors (clean maps) -----------------------------------------
\begin{frame}{Too many people for a few rooms}
\body
The building is large: two floors and ninety-eight rooms. The problem is
not the headcount but the \emp{concentration}: demand collapses onto the three masterpiece rooms
upstairs and Caravaggio downstairs, so those rooms run far over capacity while the rest sit calm.
The crowd is itself mixed, \emp{60\%} Instagram tourists, \emp{30\%} normal tourists and
\emp{10\%} art lovers.
\vfill
\begin{columns}[c]
\begin{column}{0.5\textwidth}\centering \emp{Second floor} (the three gated masterpieces)\\[3pt]
\includegraphics[width=0.94\linewidth]{\asset{uffizi_floor2.png}}\end{column}
\begin{column}{0.5\textwidth}\centering \emp{First floor} (Caravaggio, free-flow)\\[3pt]
\includegraphics[width=0.94\linewidth]{\asset{uffizi_floor1.png}}\end{column}
\end{columns}
\vfill
\centering The real Gallerie degli Uffizi plan: the famous rooms cluster off one corridor on each floor.
\end{frame}

% --- 4. Two streams ---------------------------------------------------------
\begin{frame}{Two streams, one project}
\body
We split the project into two parts that stay separate until the very end.
\vfill
\begin{columns}[T]
\begin{column}{0.5\textwidth}
\textbf{Stream 1: crowd simulator.}\\[0.5em]
Run the whole 60/30/10 population through a full day, minute by minute.
\begin{itemize}\setlength\itemsep{0.7em}
\item Compare the museum \emp{as it is} against a Pigovian redesign (booked slots, pricing, hours).
\item The \emp{director's gate}: adopt it only if welfare rises and revenue does not fall.
\item It passes: welfare \good{+31\%}, revenue preserved.
\end{itemize}
\end{column}
\begin{column}{0.5\textwidth}
\textbf{Stream 2: RL on top (our focus).}\\[0.5em]
Take that redesigned museum as fixed and ask the sharper question.
\begin{itemize}\setlength\itemsep{0.7em}
\item Not the average visitor: the \emp{optimal} one.
\item When to book, which masterpieces and how to pace the day.
\item All under genuine uncertainty about the crowd.
\end{itemize}
\end{column}
\end{columns}
\vfill
\centering \emp{The gate passes first; only then do we study the optimal visit.}
\end{frame}

% --- 5. Foundations: the toy ------------------------------------------------
\begin{frame}{Foundations: a tabular proof of concept}
\body
\begin{columns}[c]
\begin{column}{0.46\textwidth}
We first prove the idea on a \emp{12-room toy}, small enough to keep an exact table of values for
every situation. Crowds rise and fall in fixed waves; the agent is never told the schedule.
\vfill
Trained by trial and error, tabular \emp{Q-learning} discovers \emp{temporal arbitrage}: enter a
masterpiece in the lull before or after its crowd peak, not when a guidebook would send everyone.
\vfill
\emp{Honest result (5 seeds):} Q and \emp{Double-Q} come out \emp{tied} ($83.3$ vs $83.1$). On a
deterministic toy there is no overestimation bias for Double-Q to remove, so it neither helps nor
hurts, exactly as the theory predicts.
\end{column}
\begin{column}{0.54\textwidth}\centering
\includegraphics[width=\linewidth]{\fig{13_baseline_comparison.png}}\\[0.4em]
\footnotesize Learning beats every hand-written rule by a wide margin; the structured rules even go
negative because they reliably mistime the crowd.
\end{column}
\end{columns}
\vfill
\centering \emp{Timing, learned from samples alone, beats every hand-written rule.}
\end{frame}

% --- 6. Scaling + masking ---------------------------------------------------
\begin{frame}{Scaling up: deep RL and action masking}
\body
The real museum destroys the table: ninety-eight rooms, about a hundred possible moves at each
step and a state with several hundred numbers. We replace the table with a neural network and
train it with \emp{MaskablePPO}, keeping unmasked PPO and DQN as controls so any gap is about the
method, not luck.
\vfill
\begin{block}{Action masking: hand the agent the rulebook}
The museum is a graph, so from any room only a few moves are legal, yet the network has $\sim$100
outputs. We set the score of every illegal move to $-\infty$ before the softmax, so it gets zero
probability. No gradient is spent learning which moves exist; all of it goes into the decision
that matters. This is mathematically a state-dependent action set and it is the single biggest
architectural choice in the project.
\end{block}
\vfill
\centering \emp{It turns out to be decisive exactly where the decision is hardest.}
\end{frame}

% --- 7. The reframe ---------------------------------------------------------
\begin{frame}{The decisive reframe: a booking problem}
\body
We spent months on a \emp{navigation} agent before admitting the obvious: a real visitor has a map
and is guided by staff, so optimising the walk is a shortest path dressed up as RL, a trivial,
fully observed problem. The genuinely uncertain, consequential decision is the \emp{booking}.
\vfill
\begin{itemize}\setlength\itemsep{0.55em}
\item \emp{Baseline (the museum as-is):} no booking at all, just walk into the crowds and take the
      masterpieces as you find them.
\item \emp{Intervened (RAMA):} the masterpieces are gated by booked time-slots, so the decision is
      the \emp{lead time} ($1$, $7$, $21$, or $35$ days ahead), then pacing to reach each booked window.
\item A crowd-dependent \emp{fill curve}: on a busy day the good slots sell out early, so booking late
      forfeits masterpieces entirely. The discount is $\gamma=0.999$, since the payoff lands at check-in.
\end{itemize}
\vfill
\centering \emp{Choosing the right MDP, not tuning the algorithm, was more than half the project.}
\end{frame}

% --- 8. Headline result -----------------------------------------------------
\begin{frame}{Result: book early when busier and it pays}
\body
\begin{columns}[c]
\begin{column}{0.62\textwidth}\centering
\includegraphics[width=\linewidth]{\fig{21_booking_grid.png}}
\end{column}
\begin{column}{0.37\textwidth}\centering
\includegraphics[width=\linewidth]{\fig{22_book_early_curve.png}}
\end{column}
\end{columns}
\vfill
Both profiles beat the matched no-booking baseline at every crowd level (five seeds, mean$\pm$std):
\good{+20\% to +39\%} for the art lover, \good{+36\% to +38\%} for the tourist, with all three
masterpieces secured. The learned lead time \emp{rises with the crowd} (right): book-early-when-busier
is not coded in, it emerges from the reward, because booking late at a packed museum misses the rooms.
\end{frame}

% --- 9. The visit on the floor plan -----------------------------------------
\begin{frame}{What a visit looks like, on the real floor plan}
\body
\centering
\includegraphics[width=0.74\linewidth]{\fig{23_walks_art_intervened.png}}
\vfill
\raggedright
\emp{The art lover's learned visit, eight rollouts per panel, across the three crowd levels.} It is
one continuous, sensible route: enter at \code{A1}, sweep the second-floor masterpieces (gold stars)
in their booked windows, descend the staircase to Caravaggio on the first floor and leave before
closing. Every step is a one-hop move to a neighbouring room, so there are no teleports.
\end{frame}

% --- 10. Algorithm comparison -----------------------------------------------
\begin{frame}{Which algorithm wins and why masking matters}
\body
\centering
\includegraphics[width=0.62\linewidth]{\fig{15_algorithm_comparison.png}}
\vfill
\raggedright
\begin{itemize}\setlength\itemsep{0.35em}
\item \emp{MaskablePPO} is best or tied in every cell and stable across seeds.
\item \emp{Masking is decisive exactly at the hardest cell:} at the packed tourist, unmasked PPO
      \emp{collapses} to \bad{2093}, below the no-booking baseline of $2680$, while MaskablePPO holds
      \good{3685}. At easier crowds the two agree; the gap appears only where precision matters.
\item \emp{DQN is crowd-fragile and high-variance} (its baseline falls $4527\!\to\!2044$): taking a max
      over noisy crowd values overestimates, the very bias the deterministic toy did not have.
\end{itemize}
\end{frame}

% --- 11. Challenge 1 --------------------------------------------------------
\begin{frame}{The journey was the project: 1, an unlearnable reward}
\body
\sdfl
{The agent never left the museum: it would squat in one room until the day ended. The obvious fix,
zeroing all reward if it failed to exit, made the behaviour \emp{worse}, not better.}
{An all-or-nothing terminal wipe is \emp{flat} across almost every policy, so the gradient is zero.
There is a cliff at the exit and a plateau everywhere else and from the plateau nothing points
toward the door.}
{Replace the wipe with a \emp{dense} egress signal: a per-step pressure that grows as closing nears
and as you get further from an exit, a finite (not infinite) no-exit penalty and a bonus for leaving.}
{Sparse, all-or-nothing rewards are correct but often unlearnable. Shape a slope the optimiser can
actually climb.}
\end{frame}

% --- 12. Challenge 2 --------------------------------------------------------
\begin{frame}{2: the agent could not see what it had to decide on}
\body
\sdfl
{The rule ``stay until satisfied, then move on'' never stabilised; the agent over- and under-dwelt
rooms seemingly at random, even with a sensible reward.}
{The right action depended on \emp{how much value it had already extracted from this room} and that
progress was \emp{not in the observation}. It was a hidden variable, so the task was a POMDP, not the
MDP we thought we had.}
{Add a per-room \emp{appreciation-progress} value to the state, so ``have I seen enough of this
room?'' becomes a function of something the policy can actually read.}
{The state must contain everything the optimal action depends on, or you are not solving the problem
you think you are.}
\end{frame}

% --- 13. Challenge 3 --------------------------------------------------------
\begin{frame}{3: reward shaping that backfired}
\body
\sdfl
{Well-meant extra terms made things worse: a ``boredom'' penalty produced \emp{looping} and a harsh
crowd penalty produced a \emp{fake art lover} that sprinted past crowded masterpieces.}
{Every term you add is a new objective the optimiser will game (reward hacking): circling a loop beat
exiting under the boredom penalty and a crowded Botticelli scored worse than skipping it.}
{\emp{Minimal} shaping: a \emp{gentle} crowd discount so a crowded masterpiece still beats skipping it,
drop the boredom penalty entirely and let the satiation term carry ``stop lingering'' on its own.}
{Stress-test every reward term by asking what the cheapest way to collect it is and whether that is
what you actually want.}
\end{frame}

% --- 14. Challenge 4 (the best trick) ---------------------------------------
\begin{frame}{4: the do-nothing trap and the trick that cracked book-early}
\body
\sdfl
{Given a free ``decline'' button, \emp{both} profiles learned to decline all three masterpieces. And
even without it, book-early was only \emp{half} learned, with a flat spread of lead times.}
{Declining is immediately safe while booking only pays off much later, a do-nothing local optimum; and
the booking reward lands hundreds of steps away, so its gradient on the booking action is faint
(long-horizon credit assignment).}
{Treat ``always books'' as a \emp{type trait} that rules the escape hatch out and add an
\emp{immediate off-pace penalty} at booking time, a local proxy that correlates with the distant cost.}
{A small immediate signal that tracks a far-off reward is what makes a long-horizon lesson learnable.
This was the single most decisive trick.}
\end{frame}

% --- 15. Multi-seed honesty -------------------------------------------------
\begin{frame}{What five seeds revealed (and why we ran them)}
\body
A single training run can be lucky or unlucky, so a single number can mislead. We retrained every cell
under \emp{five} training seeds and report the mean and the spread across them.
\vfill
\begin{itemize}\setlength\itemsep{1.0em}
\item \emp{The tourist is rock-stable}: essentially zero variance at every crowd level.
\item \emp{The art lover degrades as the museum fills}: the gain falls from $+39\%$ to $+29\%$ to $+20\%$.
\item \emp{The honest catch:} at \emp{maximum crowd} the art lover is \emp{bimodal}, four of the five seeds
      reach $\sim$\good{+41\%} but one \bad{collapses} below baseline, so the cell averages \warn{+20\%}
      with a wide error bar. We report the spread, not the lucky seed.
\end{itemize}
\vfill
\centering \emp{Reporting variance honestly is part of the result, not a footnote.}
\end{frame}

% --- 16. Scoreboard ---------------------------------------------------------
\begin{frame}{What worked best and worst}
\body
\begin{columns}[T]
\begin{column}{0.5\textwidth}
{\large\good{\textbf{Worked best}}}
\begin{itemize}\setlength\itemsep{0.55em}
\item \emp{Action masking}, the single biggest lever.
\item \emp{Reframing} navigation into booking: the right MDP.
\item \emp{An immediate proxy} for the delayed reward.
\item \emp{Exposing appreciation progress}: Markov again.
\item \emp{Dense egress shaping}: a learnable way out.
\end{itemize}
\end{column}
\begin{column}{0.5\textwidth}
{\large\bad{\textbf{Worked worst}}}
\begin{itemize}\setlength\itemsep{0.55em}
\item \emp{All-or-nothing exit wipe}: zero gradient.
\item \emp{A free decline button}: a do-nothing trap.
\item \emp{A harsh crowd penalty}: a fake art lover.
\item \emp{A boredom penalty}: a perverse looping incentive.
\item \emp{Greedy evaluation} of a stochastic policy.
\end{itemize}
\end{column}
\end{columns}
\vfill
\centering \emp{The failures and their diagnoses, are the project; the numbers are where they stopped.}
\end{frame}

% --- 17. Conclusion + limitations -------------------------------------------
\begin{frame}{Conclusion and what is debatable}
\body
\begin{itemize}\setlength\itemsep{0.9em}
\item \emp{The result:} given the director-safe redesign, the optimal art lover and tourist
      \emp{book early, earlier when busier}, secure all three masterpieces and beat the no-booking
      baseline by \good{+20\% to +39\%} (art) and \good{+36\% to +38\%} (tourist).
\item \emp{The journey was the project:} a chain of instructive failures, each a named RL lesson,
      unlearnable reward, partial observability, reward hacking, the do-nothing trap, credit assignment.
\item \emp{Honestly debatable:} the 60/30/10 split, the dwell and crowd-density calibration and the room
      positioning are modelling choices (the floor-plan maps let you judge them); the fill curve is
      calibrated rather than fitted to booking data; and the as-is museum abstracts away the real entry
      choice (pre-book and skip the queue, or pay less and queue at the door).
\end{itemize}
\end{frame}

% --- 18. Thank you ----------------------------------------------------------
{\setbeamertemplate{footline}{}\begin{frame}[c]
\centering
{\Huge Thank you.}\\[0.9em]
{\large Questions?}\\[1.6em]
\emp{Code, figures and a plain-language guide to every algorithm:}\\[0.3em]
\code{github.com/mocl20/Uffizi}
\end{frame}}

\end{document}
